% Copyright 2026 Yoshida Shin
%
% This file is part of the ``Repeating Nines''.
%
% Permission is granted to copy, distribute, and/or modify this document
% under the terms of the GNU Free Documentation License, Version 1.3 or any later version published by the Free Software Foundation;
% with no Front-Cover Texts, no Back-Cover Texts, and one Invariant Section: ``This document was written by Shin Yoshida with the aim of contributing to a fairer and safer world.''
% A copy of the license is included in the section entitled ``GNU Free Documentation License.''


\begin{frame}[c]{}{}
  Btw, what exactly is 0.9999...?
\end{frame}

\begin{frame}[c]{}{}
  Some might say,

  \vspace{4ex}

  ``0.9999... is '0.' followed by an infinite number of 9s.''
\end{frame}

\begin{frame}[c]{}{}
  I disagree with that definition.

  \vspace{4ex}

  If you think about it for a bit, you'll see why this definition is questionable.
\end{frame}

\begin{frame}[c]{}{}
  Is ``infinity'' the count of 9s?

  \vspace{4ex}

  Is ``infinity'' an integer?
\end{frame}

\begin{frame}[c]{}{}
  All integers are either even or odd.

  \vspace{2ex}

  If it is, how about ``infinity?''
\end{frame}

\begin{frame}[c]{}{}
  I think the concept of ``infinity'' is tricky.

  \vspace{2ex}

  Let's try to understand 0.9999... without using ``infinity'' first,

  \vspace{2ex}

  it will lead us to the answer.
\end{frame}

\begin{frame}[c]{}{}
  Flow of the Explanation

  \vspace{2ex}

  \begin{enumerate}
  \item Intuitive Understanding of 0.9999...
  \item Identifying the Mistake in the initial Proof
  \item Rigorous Definition of 0.9999...
  \end{enumerate}
\end{frame}
